% GENERATED FILE. Edit canonical lessons, then run npm run generate:books.
% canonical-source-hash: fnv1a64:cf4a080e4772676c
% canonical-lessons: FR-C321-il-etait-une-fois, FR-C321-fee, FR-C321-sorciere, FR-C321-dragon, FR-C321-fable

\chapter{Il était une fois}
\label{ch:fr-il-etait-une-fois}
\hlchaptermodality{321}

\begin{chapteropening}
\textbf{By the end of this chapter:} \emph{I can start a fairy tale and name who lives in it: once upon a time, a fairy, a witch, a dragon, and a fable.}

Complete the last lesson of chapter 321: I can start a fairy tale and name who lives in it: once upon a time, a fairy, a witch, a dragon, and a fable.
\end{chapteropening}

\section[il était une fois]{il était une fois — once upon a time}

\label{lesson:FR-C321-il-etait-une-fois}



\begin{quote}
Before the new one: say the French for the candle, then the French for the wish.
\end{quote}

\subsection*{You'll want to know: il était une fois}
\textbf{il était une fois} — \textquotedblleft{}once upon a time\textquotedblright{}.

\begin{quote}
Il était une fois une petite fille qui vivait dans la forêt.
\end{quote}

\begin{quote}
\emph{Once upon a time there was a little girl who lived in the forest.}
\end{quote}

The opening of most fairy tales. Literally 'there was one time'.

\begin{cousinweb}
\emph{fois} is inherited from Latin \emph{vicēs} \textquotedblleft{}turns, times\textquotedblright{} (Lewis \& Short); \emph{était} is a form of \emph{être}, whose imperfect is inherited from Latin \emph{stāre} \textquotedblleft{}to stand\textquotedblright{} (EtymDB).
\end{cousinweb}

\begin{grammarlens}[title={What you've built}]
One more word for telling, explaining and describing at B1.
\end{grammarlens}

\subsection*{Your turn}
\begin{itemize}
\raggedright
  \item \emph{Say it:} \emph{il était une fois}
  \item \emph{Say it:} \emph{il était une fois}, once more
  \item \emph{Say it:} a sentence of your own with \emph{il était une fois}
  \item \emph{Recall:} say the French for the candle, then the French for the wish, then say \emph{il était une fois} again
\end{itemize}

\begin{tcolorbox}[breakable,colback=teal!4,colframe=teal!35!black,title={Before you move on}]
What does il était une fois mean? (\textquotedblleft{}Once upon a time\textquotedblright{}.) Say it once more.
\end{tcolorbox}
\section[la fée]{la fée — the fairy}

\label{lesson:FR-C321-fee}



\begin{quote}
Before the new one: say the French for the wish, then the French for once upon a time.
\end{quote}

\subsection*{You'll want to know: la fée}
\textbf{la fée} — \textquotedblleft{}the fairy\textquotedblright{}.

\begin{quote}
La bonne fée a aidé la jeune fille.
\end{quote}

\begin{quote}
\emph{The good fairy helped the girl.}
\end{quote}

\emph{Une bonne fée} — a fairy godmother.

\begin{cousinweb}
\emph{fée} is inherited from Vulgar Latin \emph{Fāta} \textquotedblleft{}goddess of fate\textquotedblright{} (EtymDB), built on Latin \emph{fātum} \textquotedblleft{}destiny\textquotedblright{}.
\end{cousinweb}

\begin{grammarlens}[title={What you've built}]
One more word for telling, explaining and describing at B1.
\end{grammarlens}

\subsection*{Your turn}
\begin{itemize}
\raggedright
  \item \emph{Say it:} \emph{la fée}
  \item \emph{Say it:} \emph{la fée}, once more
  \item \emph{Say it:} a sentence of your own with \emph{la fée}
  \item \emph{Recall:} say the French for the wish, then the French for once upon a time, then say \emph{la fée} again
\end{itemize}

\begin{tcolorbox}[breakable,colback=teal!4,colframe=teal!35!black,title={Before you move on}]
What does la fée mean? (\textquotedblleft{}The fairy\textquotedblright{}.) Say it once more.
\end{tcolorbox}
\section[la sorcière]{la sorcière — the witch}

\label{lesson:FR-C321-sorciere}



\begin{quote}
Before the new one: say the French for once upon a time, then the French for the fairy.
\end{quote}

\subsection*{You'll want to know: la sorcière}
\textbf{la sorcière} — \textquotedblleft{}the witch\textquotedblright{}.

\begin{quote}
La sorcière vivait seule dans une maison noire.
\end{quote}

\begin{quote}
\emph{The witch lived alone in a black house.}
\end{quote}

The masculine \emph{le sorcier} is a wizard.

\begin{cousinweb}
\emph{sorcière} is the feminine of \emph{sorcier}, inherited from Vulgar Latin \emph{sortiārius} \textquotedblleft{}one who casts lots\textquotedblright{}, built on \emph{sors} \textquotedblleft{}a lot, fate\textquotedblright{} (Lewis \& Short).
\end{cousinweb}

\begin{grammarlens}[title={What you've built}]
One more word for telling, explaining and describing at B1.
\end{grammarlens}

\subsection*{Your turn}
\begin{itemize}
\raggedright
  \item \emph{Say it:} \emph{la sorcière}
  \item \emph{Say it:} \emph{la sorcière}, once more
  \item \emph{Say it:} a sentence of your own with \emph{la sorcière}
  \item \emph{Recall:} say the French for once upon a time, then the French for the fairy, then say \emph{la sorcière} again
\end{itemize}

\begin{tcolorbox}[breakable,colback=teal!4,colframe=teal!35!black,title={Before you move on}]
What does la sorcière mean? (\textquotedblleft{}The witch\textquotedblright{}.) Say it once more.
\end{tcolorbox}
\section[le dragon]{le dragon — the dragon}

\label{lesson:FR-C321-dragon}



\begin{quote}
Before the new one: say the French for the fairy, then the French for the witch.
\end{quote}

\subsection*{You'll want to know: le dragon}
\textbf{le dragon} — \textquotedblleft{}the dragon\textquotedblright{}.

\begin{quote}
Le dragon dormait sur une montagne d'or.
\end{quote}

\begin{quote}
\emph{The dragon was sleeping on a mountain of gold.}
\end{quote}

In old armies, \emph{un dragon} was also a horse soldier.

\begin{cousinweb}
\emph{dragon} is inherited from Latin \emph{dracō}, \emph{dracōnem} (EtymDB), a borrowing from Greek \emph{drákōn} \textquotedblleft{}serpent\textquotedblright{}.
\end{cousinweb}

\begin{grammarlens}[title={What you've built}]
One more word for telling, explaining and describing at B1.
\end{grammarlens}

\subsection*{Your turn}
\begin{itemize}
\raggedright
  \item \emph{Say it:} \emph{le dragon}
  \item \emph{Say it:} \emph{le dragon}, once more
  \item \emph{Say it:} a sentence of your own with \emph{le dragon}
  \item \emph{Recall:} say the French for the fairy, then the French for the witch, then say \emph{le dragon} again
\end{itemize}

\begin{tcolorbox}[breakable,colback=teal!4,colframe=teal!35!black,title={Before you move on}]
What does le dragon mean? (\textquotedblleft{}The dragon\textquotedblright{}.) Say it once more.
\end{tcolorbox}
\section[la fable]{la fable — the fable}

\label{lesson:FR-C321-fable}



\begin{quote}
Before the new one: say the French for the witch, then the French for the dragon.
\end{quote}

\subsection*{You'll want to know: la fable}
\textbf{la fable} — \textquotedblleft{}the fable\textquotedblright{}.

\begin{quote}
Les enfants apprennent des fables à l'école.
\end{quote}

\begin{quote}
\emph{Children learn fables at school.}
\end{quote}

French children learn fables by heart, with animals that talk.

\begin{cousinweb}
\emph{fable} comes from Old French \emph{fable}, a learned borrowing from Latin \emph{fābula} \textquotedblleft{}a story\textquotedblright{} (EtymDB), built on \emph{fārī} \textquotedblleft{}to speak\textquotedblright{}.
\end{cousinweb}

\begin{grammarlens}[title={What you've built}]
One more word for telling, explaining and describing at B1.
\end{grammarlens}

\subsection*{Your turn}
\begin{itemize}
\raggedright
  \item \emph{Say it:} \emph{la fable}
  \item \emph{Say it:} \emph{la fable}, once more
  \item \emph{Say it:} a sentence of your own with \emph{la fable}
  \item \emph{Recall:} say the French for the witch, then the French for the dragon, then say \emph{la fable} again
\end{itemize}

\begin{tcolorbox}[breakable,colback=teal!4,colframe=teal!35!black,title={Before you move on}]
What does la fable mean? (\textquotedblleft{}The fable\textquotedblright{}.) Say it once more.
\end{tcolorbox}
